```latex
\documentclass[11pt]{article}
\usepackage{amsmath,amssymb,amsthm}
\usepackage{mathtools}
\usepackage{bbm}

\newtheorem{theorem}{Theorem}[section]
\newtheorem{lemma}[theorem]{Lemma}
\newtheorem{proposition}[theorem]{Proposition}
\newtheorem{corollary}[theorem]{Corollary}
\theoremstyle{remark}
\newtheorem{remark}[theorem]{Remark}

\DeclareMathOperator{\div}{div}
\DeclareMathOperator{\E}{\mathbb{E}}
\DeclareMathOperator{\Var}{Var}
\DeclareMathOperator{\tr}{tr}
\newcommand{\R}{\mathbb{R}}
\newcommand{\norm}[1]{\left\|#1\right\|}
\newcommand{\abs}[1]{\left|#1\right|}
\newcommand{\one}{\mathbf{1}}

\title{Minimax rate for a non-linear Darcy inverse problem}
\author{}
\date{}

\begin{document}
\maketitle

\section{Setup and main result}

Let $D\subset \R^d$, $d\ge 2$, be the Euclidean unit ball. For $f\in C^\alpha(D)$,
$\alpha>1$, with $f\ge c_0>0$, let $u=u_f$ be the unique solution of
\[
\div(f\nabla u)=1\quad \text{in }D,\qquad u|_{\partial D}=0.
\]
We write
\[
F(f)=u_f.
\]
For small fixed $r\in(0,1)$, set
\[
\mathcal F(\alpha,r)=\{f\in C^\alpha(D):\ \norm{f-\one}_{C^\alpha(D)}\le r\}.
\]
We observe
\[
Y_n=u_f+\frac1{\sqrt n}\mathbb W,
\]
where $\mathbb W$ is Gaussian white noise on $L^2(D)$.  The minimax risk is
\[
R_n(\alpha,r)=\inf_{\hat f_n}\sup_{f_0\in\mathcal F(\alpha,r)}
\E_{f_0}\bigl[\norm{\hat f_n-f_0}_{L^2(D)}^2\bigr].
\]

\begin{theorem}\label{thm:main}
For every sufficiently small fixed $r\in(0,1)$ and every $\alpha>1$,
\[
R_n(\alpha,r)\asymp n^{-\frac{2\alpha}{2\alpha+d+3}}
\]
up to logarithmic factors in $n$.
\end{theorem}

The remainder of this note proves Theorem~\ref{thm:main}. The main point is that the
linearization of the forward map at $f=\one$ is an anisotropic smoothing operator:
it smooths radial oscillations by one derivative, but angular oscillations by two
derivatives. The latter dominate the minimax rate.

\section{Linearization and singular value asymptotics}

The background solution for $f=\one$ is
\[
u_0(x)=\frac{|x|^2-1}{2d},
\]
which satisfies $\Delta u_0=1$ and $u_0|_{\partial D}=0$. Write
\[
f=1+g,\qquad \norm{g}_{C^\alpha(D)}\le r.
\]
Then
\[
u_f=u_0+v,
\]
where $v$ solves
\[
\Delta v+\div(g\nabla u_0)+\div(g\nabla v)=0,\qquad v|_{\partial D}=0.
\]
Thus the Fr\'echet derivative of $F$ at $\one$ is
\[
Tg=v_1,
\]
where
\begin{equation}\label{eq:T}
\Delta v_1=-\div(g\nabla u_0),\qquad v_1|_{\partial D}=0.
\end{equation}
Since $\nabla u_0=x/d$, in spherical coordinates $x=\rho\theta$, $\rho\in(0,1)$,
$\theta\in S^{d-1}$,
\[
\div(g\nabla u_0)=\frac1d\div(gx)
=\frac1d\bigl(\rho\partial_\rho g+d\,g\bigr).
\]
Thus the angular variables enter only through the Laplacian inverse.

\begin{lemma}\label{lem:sv}
Let $\sigma_{k,l}$ denote the singular values of $T$ corresponding to radial
frequency $k$ and angular frequency $l$. Then, for $k,l\ge1$,
\[
\sigma_{k,l}\asymp \frac{k}{k^2+l^2}.
\]
In particular, for angular frequencies $l\asymp L$ and bounded radial frequency
$k\asymp 1$, we have $\sigma_{k,l}\asymp L^{-2}$, and the multiplicity of such
modes is $\asymp L^{d-1}$.
\end{lemma}

\begin{proof}[Proof sketch]
Separating variables in spherical harmonics on $S^{d-1}$, the Laplacian becomes
\[
\Delta \;\longleftrightarrow\;
\rho^{-(d-1)}\partial_\rho(\rho^{d-1}\partial_\rho)
-\frac{l(l+d-2)}{\rho^2}.
\]
For large $l$ and radial frequency $k$, the inverse Laplacian has symbol
$(k^2+l^2)^{-1}$. The operator $g\mapsto \div(g\nabla u_0)$ is a first-order
radial operator with symbol of size $k$, independent of $l$. Therefore
\[
\sigma_{k,l}\asymp \frac{k}{k^2+l^2}.
\]
The dimension of the space of spherical harmonics of degree $l$ is
$\asymp l^{d-2}$, so the number of angular modes up to $L$ is $\asymp L^{d-1}$.
\end{proof}

\section{Lower bound}

We use Fano's method. The construction is based on angular wavepackets of
bandwidth $L$; these are the modes for which $T$ has singular values of order
$L^{-2}$ and which have multiplicity $L^{d-1}$.

Let $\{\psi_j\}_{j=1}^{M}$, $M\asymp L^{d-1}$, be an orthonormal family of
angular wavepackets on $S^{d-1}$ of bandwidth $L$, supported on disjoint caps of
radius $\sim L^{-1}$. Such a system satisfies
\[
\norm{\psi_j}_{L^2(S^{d-1})}=1,\qquad
\norm{\psi_j}_{C^\alpha(S^{d-1})}\le C L^{\alpha+\frac{d-1}{2}}.
\]
Let $\phi\in C_c^\infty((1/4,3/4))$ be a radial cutoff with
$\norm{\phi}_{L^2(\rho^{d-1}d\rho)}=1$. For
\[
a_L:=c r L^{-\alpha-\frac{d-1}{2}},
\]
define
\[
g_j(x):=a_L \phi(\rho)\psi_j(\theta),\qquad j=1,\dots,M.
\]
Then, if $c>0$ is small enough,
\[
\norm{g_j}_{C^\alpha(D)}\le r.
\]
For $\theta=(\theta_1,\dots,\theta_M)\in\{\pm1\}^M$, set
\[
g_\theta:=\sum_{j=1}^M \theta_j g_j,\qquad
f_\theta:=1+g_\theta.
\]
Since the angular supports are disjoint, the $C^\alpha$-norm of $g_\theta$ is
comparable to the maximum of the norms of the individual bumps, hence
\[
f_\theta\in\mathcal F(\alpha,r)
\]
for all $\theta$, after possibly reducing the constant $c$.

For two sign vectors $\theta,\theta'$, let
$H(\theta,\theta')$ be their Hamming distance. Then
\[
\norm{g_\theta-g_{\theta'}}_{L^2(D)}^2
=4 H(\theta,\theta')\, a_L^2.
\]
By the Gilbert--Varshamov bound, there exists a subset
$\Theta\subset\{\pm1\}^M$ of cardinality at least $e^{c_0 M}$ such that
\[
H(\theta,\theta')\ge \frac M4\qquad \forall \theta\neq\theta'\in\Theta.
\]
Thus, for $\theta\neq\theta'$,
\begin{equation}\label{eq:delta}
\norm{g_\theta-g_{\theta'}}_{L^2(D)}^2
\ge c M a_L^2
\asymp L^{d-1}\cdot L^{-2\alpha-(d-1)}
=L^{-2\alpha}.
\end{equation}

It remains to control the distance between the corresponding forward solutions.
By Lemma~\ref{lem:sv}, for a function $w$ built from angular wavepackets of
bandwidth $L$ and bounded radial frequency,
\[
\norm{Tw}_{L^2(D)}\le C L^{-2}\norm{w}_{L^2(D)}.
\]
Moreover, the non-linear remainder is of lower order. Indeed, subtracting the
PDEs for two conductivities $1+g$ and $1+h$ and linearizing at $f=\one$ gives
\[
F(1+g)-F(1+h)=T(g-h)+R(g,h),
\]
with
\[
\norm{R(g,h)}_{L^2(D)}\le C r \norm{T(g-h)}_{L^2(D)}
\]
for all $g,h$ in the $C^\alpha$-ball of radius $r$; this follows from the usual
Schauder estimates and the fact that the non-linearity is quadratic. Hence, for
$\theta,\theta'\in\Theta$,
\[
\norm{u_{f_\theta}-u_{f_{\theta'}}}_{L^2(D)}
\le C L^{-2}\norm{g_\theta-g_{\theta'}}_{L^2(D)}
\le C L^{-\alpha-2}.
\]
Consequently the squared $L^2$-distance between the corresponding signals is at most
\[
\varepsilon_L^2:=C L^{-2\alpha-4}.
\]

The observations for two different sign vectors are Gaussian with the same
covariance and means separated by at most $\varepsilon_L$ in $L^2$. Therefore
the Kullback--Leibler divergence between the two distributions is bounded by
\[
\frac n2 \varepsilon_L^2\le C n L^{-2\alpha-4}.
\]
Fano's inequality applies if this is at most $c\log|\Theta|\asymp M\asymp L^{d-1}$.
Thus we choose
\[
L\asymp n^{\frac1{2\alpha+d+3}}.
\]
By \eqref{eq:delta}, the resulting $\theta$-dependent conductivities are
pairwise separated in $L^2$ by at least
\[
c L^{-\alpha}
\asymp n^{-\frac{\alpha}{2\alpha+d+3}}.
\]
Fano's lemma yields
\[
R_n(\alpha,r)\ge c n^{-\frac{2\alpha}{2\alpha+d+3}}.
\]
This proves the lower bound.

\section{Upper bound}

For the upper bound we use a standard non-linear sieve estimator. Let
$\mathcal V_L$ be a finite-dimensional space spanned by angular wavepackets of
bandwidth at most $L$ and radial basis functions of frequency at most $L$. Its
dimension is $N_L\asymp L^d$. It has the usual approximation property: for every
$g\in C^\alpha(D)$,
\[
\inf_{\bar g\in\mathcal V_L}\norm{g-\bar g}_{L^2(D)}
\le C L^{-\alpha}\norm{g}_{C^\alpha(D)}.
\]
Define
\[
\mathcal G_L:=\{g\in\mathcal V_L:\ \norm{g}_{C^\alpha(D)}\le r\}.
\]
Let $\hat g_L$ be the constrained non-linear least squares estimator
\[
\hat g_L:=
\arg\min_{g\in\mathcal G_L}
\norm{Y_n-u_0-\bigl(F(1+g)-u_0\bigr)}_{L^2(D)}^2.
\]
Set $\hat f_n:=1+\hat g_L$.

The key stability estimate is the following. For $g,h\in\mathcal G_L$, let
$p=F(1+g)-F(1+h)$. Subtracting the two PDEs gives
\[
\div((1+g)\nabla p)=-\div((g-h)\nabla u_h),
\]
where $u_h=F(1+h)$. Since $u_h$ is a small $C^\alpha$-perturbation of $u_0$,
the vector field $\nabla u_h$ is a small perturbation of $x/d$. Solving the
transport equation
\[
\div(w\nabla u_h)=-h
\]
for $w=g-h$ gives
\[
\norm{g-h}_{L^2(D)}\le C\norm{p}_{H^2(D)}.
\]
On the other hand, because $g,h\in\mathcal V_L$, elliptic regularity gives the
Bernstein-type inequality
\[
\norm{p}_{H^2(D)}\le C L^2\norm{p}_{L^2(D)}.
\]
Therefore the restricted forward map is bi-Lipschitz:
\begin{equation}\label{eq:bilip}
c L^{-2}\norm{g-h}_{L^2(D)}
\le \norm{F(1+g)-F(1+h)}_{L^2(D)}
\le C\norm{g-h}_{L^2(D)}.
\end{equation}
The lower bound in \eqref{eq:bilip} is the only stability fact needed.

Using \eqref{eq:bilip}, the standard contraction inequality for non-linear least
squares over a finite-dimensional compact parameter space gives, for every
$\bar g\in\mathcal G_L$,
\[
\E_{f_0}\norm{\hat g_L-g_0}_{L^2(D)}^2
\le C\norm{g_0-\bar g}_{L^2(D)}^2
+C n^{-1}\sum_{\lambda\in\Lambda_L}\sigma_\lambda^{-2},
\]
where $\sigma_\lambda$ are the singular values of the linearized operator $T$
restricted to $\mathcal V_L$, and $\Lambda_L$ indexes the corresponding modes.
The first term is the approximation bias, bounded by $C L^{-2\alpha}$. The
second term is the variance. From Lemma~\ref{lem:sv},
\[
\sum_{\lambda\in\Lambda_L}\sigma_\lambda^{-2}
\asymp \sum_{l\le L} l^{d-2} l^4
\asymp L^{d+3}.
\]
Hence
\[
\E_{f_0}\norm{\hat f_n-f_0}_{L^2(D)}^2
\le C\bigl(L^{-2\alpha}+n^{-1}L^{d+3}\bigr).
\]
Choosing
\[
L\asymp n^{\frac1{2\alpha+d+3}}
\]
balances the two terms and gives
\[
\E_{f_0}\norm{\hat f_n-f_0}_{L^2(D)}^2
\le C n^{-\frac{2\alpha}{2\alpha+d+3}},
\]
up to logarithmic factors. This proves the upper bound.

\section{Conclusion}

Combining the lower and upper bounds proves Theorem~\ref{thm:main}: for small
fixed $r$, the minimax rate of estimation of $f$ in $L^2$-mean squared error is
\[
\boxed{\ n^{-\frac{2\alpha}{2\alpha+d+3}}\ }
\]
up to logarithmic factors.

\end{document}
```